% tab_efeito_antipadrao; regime quente; experimentos [1, 2, 4, 5, 6, 7, 8, 9, 10, 11, 12, 13, 14, 15, 16, 17, 18, 19, 20]; criterios_sha 23e5e61d4e23e46b
% requer \usepackage{booktabs}; arquivo gerado por tcclab.relatorio (nao editar a mao)
{\small\setlength{\tabcolsep}{4pt}
\begin{tabular}{lrlrrrrl}
\toprule
Estrato & $n$ & Métrica & Redução mediana & IC95\% (bootstrap) & $p$ & $p$ (Holm) &  \\
\midrule
Todas & 17 & Leituras lógicas & 99,9\% & [99,0; 100,0] & $<$0,001 & $<$0,001 &  \\
Todas & 17 & CPU (ms) & 99,0\% & [96,9; 99,8] & $<$0,001 & $<$0,001 &  \\
Todas & 17 & Decorrido (ms) & 99,0\% & [97,0; 99,9] & $<$0,001 & $<$0,001 &  \\
\midrule
AP1 & 4 & Leituras lógicas & 99,9\% & [99,4; 100,0] & 0,125 & 0,375 & $^{\dagger}$ \\
AP1 & 4 & CPU (ms) & 99,4\% & [98,3; 100,0] & 0,125 & 0,375 & $^{\dagger}$ \\
AP1 & 4 & Decorrido (ms) & 99,4\% & [98,5; 100,0] & 0,125 & 0,375 & $^{\dagger}$ \\
\midrule
AP2 & 4 & Leituras lógicas & 63,1\% & [6,4; 99,0] & 0,125 & 0,375 & $^{\dagger}$ \\
AP2 & 4 & CPU (ms) & 52,2\% & [5,6; 97,2] & 0,125 & 0,375 & $^{\dagger}$ \\
AP2 & 4 & Decorrido (ms) & 51,6\% & [7,2; 97,6] & 0,125 & 0,375 & $^{\dagger}$ \\
\midrule
AP3 & 4 & Leituras lógicas & 100,0\% & [99,6; 100,0] & 0,125 & 0,375 & $^{\dagger}$ \\
AP3 & 4 & CPU (ms) & 99,6\% & [96,9; 100,0] & 0,125 & 0,375 & $^{\dagger}$ \\
AP3 & 4 & Decorrido (ms) & 99,6\% & [97,0; 100,0] & 0,125 & 0,375 & $^{\dagger}$ \\
\midrule
AP4 & 5 & Leituras lógicas & 100,0\% & [98,9; 100,0] & 0,062 & 0,188 & $^{\dagger}$ \\
AP4 & 5 & CPU (ms) & 99,8\% & [91,3; 100,0] & 0,062 & 0,188 & $^{\dagger}$ \\
AP4 & 5 & Decorrido (ms) & 99,8\% & [87,0; 100,0] & 0,062 & 0,188 & $^{\dagger}$ \\
\bottomrule
\end{tabular}\par
\vspace{3pt}
\begin{minipage}{\linewidth}\footnotesize
Regime quente. Redução mediana = (1 $-$ mediana das razões depois/antes entre consultas) $\times$ 100; $p$: Wilcoxon de postos sinalizados sobre o log da razão (bilateral), Holm sobre 3 métricas, $\alpha$ = 0,05. Sem correção entre estratos.\par
$^{\dagger}$ Exploratório: menos de 6 consultas ou $p$ mínimo atingível acima de $\alpha$ (o teste não pode rejeitar $H_0$).
\end{minipage}
}
